\paragraph{Compare with model checking (Reviewer B and E)} It is
not oblivious how to preform "model checking for PAT" (Reviewer
E). First of all, to enable model check on open systems, users still
need to close the system by providing the input generation program
that interacts with handlers (actors). Moreover, \PAT provides the
\emph{local} specification for each handlers, it is non-trivial how to
lift them into \emph{global} level since these specifications may
conflict with each others. Synthesizing controller by naive
enumeration is exactly an possible implementation of this lifting
procedure. Finally, this model checking still cause nontrivial vast
search space and requires smart searching strategies.

Given concerns above, instead of "fully explorable" (Reviewer E), \PAT
offers an incomplete but a light way approach that efficiently
explores executions guided by the properties we aim to reason
about. Thus, \PAT doesn't work as "an abstract version of the
implementation" (Reviewer C) or "a complete, symbolic description of
the system under test" (Reviewer E) in model check setting, but more
like specifications that are sufficient to help us prune out
executions that cannot violate the given properties in the test
setting.

\paragraph{Expressivity and Limitations of \PAT (All Reviewers)} As mentioned by
Reviewer E, the ghost variables increase the expressive power of our
SFAs to be "more similar to, e.g., data/register automata that can
save values". We would also like to highlight that \PAT provide
constraints over each actor on its \emph{local} view where the
controller has an \emph{global} view. This design make the trace in
our approach is not just "one global trace of all events in the
system" (Reviewer E), but an merged views of all actors. There can be
multiple events for the same operation, comes from different actors
under different local view. Then the global property can specify these
events are consistent or not. The "multiple traces per thread"
properties (e.g., sequential consistency) and network behaviors can
then expressed in this way, where we provides examples in the details
response to Reviewer E. We also provide the \PAT and explanation for
all benchmarks in the attached file of author response.

The limitations of \PAT mostly stem from the limitations of SFAs. For
example, SFAs cannot express counting properties (e.g., an event
appearing twice as often as another event), as this exceeds the
expressive power of regular languages. Additionally, since we only
consider finite traces, we cannot express properties involving
infinite traces (e.g., an event appearing finitely many times in the
future). The "liveness property" mentioned in our paper actually
refers to "eventually exists within a finite number of steps,"
following an informal interpretation of liveness: "something good will
eventually happen." We apologize for this misuse and promise to
clarify it in our paper revision.

\paragraph{LTL$_f$, SFA, and Abstract Trace (All reviewers)} We clarify that {\sf Clouseau} indeed is based on SFAs, and it can accepts any frontend language (e.g., symbolic regex commended by reviewers) that can be compiled into SFAs. Then our abstract trace can be treat as "symbolic regex" without \emph{top-level} union $\Pi ::= \msgB{op}{\overline{x}}{\phi} ~|~ \Pi \seqA \Pi ~|~ A^*$. Then, we can normalize each SFA into a finite set of abstract traces, as normalization process will preserve the star term instead of unfolding it. Current presentation in our paper attempts to mimic this simple idea in LTL$_f$, which has led to a lot of confusion. We promise to present the abstract trace and normalization in symbolic regex in our revised paper.



\section{Attached File}

\paragraph{Detailed Response}

\paragraph{Expressivity and Limitations of \PAT (All Reviewers)}

We clarify that the Prophecy Automata Type (\PAT) \(\rg{H}{A}{F}\) indeed uses SFAs (i.e., \(H\), \(A\), and \(F\)) to encode history, current, and future traces. Similar to other trace-based refinement types \cite{HAT}, LTL\(_f\) serves as a user-friendly frontend language for SFAs, improving readability, and is later compiled into SFAs in {\sf Clouseau}. We also allow users to write symbolic regular language (SRL) or linear dynamic logic (LDL\(_f\)) \cite{LTLf}, which are more powerful than LTL\(_f\), as long as they can be compiled into SFAs.

The limitations of \PAT mostly stem from the limitations of SFAs. For example, SFAs cannot express counting properties (e.g., an event appearing twice as often as another event), as this exceeds the expressive power of regular languages. Additionally, since we only consider finite traces, we cannot express properties involving infinite traces (e.g., an event appearing finitely many times in the future). The "liveness property" mentioned in our paper actually refers to "eventually exists within a finite number of steps," following an informal interpretation of liveness: "something good will eventually happen." We apologize for this misuse and promise to clarify it in our paper revision.

Despite these limitations inherited from SFAs, our approach can still express interesting properties with the help of ghost variables, which connect events within SFAs, and prophecy automata, which allow to express \emph{local} views for each handler (actor). All verification conditions in our setting check whether the qualifiers in symbolic events are satisfiable to prevent unrealizable traces (Reviewer A). Thus, adding ghost variables (universally quantified variables) still keeps our VCs within EPR, as long as the qualifiers remain in EPR (thanks to Reviewer E; we will highlight this point in the paper revision).  
We address the expressivity concerns raised by Reviewer E with the following examples.
\begin{enumerate}
    \item Expressing "exactly one". We can add a request ID as the \(\I{rid}\) field for request and response events, then constrain it so that there is at most one request and one response for each ghost variable \(i\). For example,
    \begin{align*}
        i{:}\Code{rId} \sarr \rg{H}{\msg{putRsp}{rid = i}}{F \land \neg \finalA (\msg{putRsp}{rid = i} \land \nextA \finalA \msg{putRsp}{rid = i})}
    \end{align*}
    \item Expressing sequential consistency, linearizability, deadlock freedom, etc. \PAT{} indeed uses a global trace but can still express properties involving multiple participants through the local views provided by \PAT{}. We demonstrate how to encode Lamport's sequential consistency property \cite{Lamport1979} with a simple example where multiple nodes (processes) perform read and write operations on a database. Each node communicates with the database using \(\eff{readReq/Resp}\) and \(\eff{writeReq/Resp}\) events, which represent the local view of each process. Meanwhile, the controller also sends auxiliary events, \(\eff{gReadReq/Resp}\) and \(\eff{gWriteReq/Resp}\), to the database to indicate a global sequential order. By default, the order of global read/write events and actual read/write events is completely independent. Sequential consistency requires that the global order matches the local view, which can be expressed by the following property:
    \begin{align*}
        \forall i{:}\Code{requestId}, &(\msg{writeReq}{\I{rid} = i} \implies \nextA \msg{gWriteReq}{\I{rid} = i}) \land 
        \\&(\msg{writeResp}{\I{rid} = i} \implies \nextA \msg{gWriteResp}{\I{rid} = i}) \land ...
    \end{align*}
    where we require that local read and write operations align with the global ones. Additionally, sequential consistency mandates that the execution remains the same as if all operations were performed in the global order. We can formalize this constraint using the \PAT{} of the database and nodes, ensuring that the database returns the last \emph{globally} written value as its read result
    {\small\begin{align*}
        i{:}\Code{rId} \garr \Code{x}\hasnty{int} \sarr \rg{\finalA \msg{gWriteReq}{v = \Code{x}} \land \nextA \neg \finalA \eff{gWriteReq}}{\lastA \msg{readReq}{\I{id} = i}}{\msg{readRsp}{v = \Code{x} \land \I{id} = i}}
    \end{align*}}
    and node expected the last \emph{local} written value as read result.
    {\small\begin{align*}
        i{:}\Code{rId} \sarr \Code{x}\hasnty{int} \sarr \rg{\finalA \msg{writeReq}{v = \Code{x}} \land \nextA \neg \eff{writeReq} \untilA \msg{readReq}{\I{id} = i} }{\lastA \msg{readRsq}{\I{id} = i}}{\allTL}
    \end{align*}}
    The \PAT{} of the database uses prophecy automata to \emph{predict} that the event \(\eff{readRsp}\) contains the last written value in the global view, while the \PAT{} of the node uses current automata to validate that the same event \(\eff{readRsp}\) returns the last written value in the local view. In general, \PAT{} allows handlers to impose multiple constraints on the same event from their respective views, effectively expressing consistency properties involving multiple participants.
    
    \item Expressing Network Behavior. In-order and out-of-order delivery can be specified using the prophecy automata in the \PAT{} of each handler. Suppose a handler sends event \(\eff{E_1}\) followed by \(\eff{E_2}\). The following \PAT{} enforces in-order delivery by ensuring that \(\eff{E_1}\) is received before \(\eff{E_2}\):
    \begin{align*}
        \rg{H}{A}{\finalA (\eff{E_1} \land \nextA \finalA \eff{E_2})}
    \end{align*}
    then the following \PAT allows that $\eff{E_1}$ and $\eff{E_2}$ appears in arbitrary order:
    \begin{align*}
        \rg{H}{A}{(\finalA \eff{E_1}) \land (\finalA \eff{E_2})}
    \end{align*}
    Node failure can be represented as a special event, \(\eff{crash}\), sent by the controller. When a node receives this event, it will no longer respond to any incoming messages. This means that an handler can handle messages normally if and only if it has not received the \(\eff{crash}\) event:
    \begin{align*}
        \rg{H \land \neg (\finalA \eff{crash})}{A}{F}
    \end{align*}
\end{enumerate}

\textbf{SJ start}

\section{Main Response}

We thank the reviewers for their detailed comments and suggestions.

\subsection{Overview}

We begin by clarifying concerns that were shared by multiple
reviewers. Specifically, we (a) characterize the expressivity and limitations of our
\PAT-based type abstraction; (b) relateg the novelty of our
methodology to other PBT and model-checking techniques; and (c)
justify our use of PATs as being particularly well-suited to
property-based testing of distributed system models. \BD{We do not
  specifically target the last point.}

We then present a detailed changelist that proposes to (a) better
clarify our methodology and its applicability to distributed systems
testing; (b) incorporate additional related work as suggested by the
reviewers; and (c) elaborate and expand our evaluation by providing
additional experimental details as well as additional experiments
demonstrating that our technique applies to test well-understood
consistency and network delivery properties.

Finally, we provide detailed responses to the questions raised by
individual reviewers.

\subsection{Expressivity, Framework, and Methodology (All Reviewers)}

\subsubsection{Expressivity}

% Point 1: PATs allow us to capture fine-grained dependencies between
% actors in a flexible way.

Beyond the ability to express temporal modalities on events, the PAT
specification language allows the use of ghost variables to capture
data dependencies among specification events, resulting in a level of
expressive power akin to data/register automata that are equipped with
memory (Reviewer E), while still being amenable to efficient SMT
encodings. This additional power allows us to capture fine-grained
dependencies between the actors in the SUT, while at the same time not
overly constraining the behavior of the complete system. These
behaviors are delegated to the synthesized controllers, allowing us to
capture and test for wide range of consistency, (e.g., XXX) and
network delivery properties (e.g., YYY) relevant to distributed
systems (Reviewers E); we elaborate on this point in the individual
responses, and provide detailed examples of these in the attached
file.

Our PAT-based specifications inherit the limitations of the SFAs that
they compile into. As one example, SFAs (even when equipped with ghost
variables) cannot express counting properties (e.g., an event
appearing twice as often as another event), as this exceeds the
expressive power of regular languages. Additionally, since we only
consider finite traces, we cannot express properties involving
infinite traces (e.g., an event appearing infinitely many times in the
future). For properties that are amenable to PBT-style automated
testing, however, our SFA representations appears to be particularly
well-suited and effective.

% Completeness:

%- l823: Our implementation resolves this nondeterminism
% in the algorithm via an efficient backtracking search procedure that
% takes the union of all successful runs in order to capture different
% orderings.
%
% - l921: Each controller synthesized by Algo rithm 1 fixes a
% particular message order for generable (i.e., environment) messages,
% so Clouseau systematically explores the space of alternative
% orderings, returning the set of all controllers found within a
% user-provided time bound.

% Unlike traditional PBT approaches that only specify a property of
% interest that should be falsified, our technique additionally allows
% rich specifications in the form of PATs (Prophecy Automata Types) to
% be attached to the actors comprising a SUT.

\subsection{Comparison to model-checking and related techniques}
% Point 2: PATs are an approximation of the behaviors of an actor, and
% do not provide a complete view (so we would still have to test).
% Even if they did, the space of possible executions is so large that
% we would have to apply all your favorite MC techniques, partial
% order reduction, etc. to keep things tractible.

There are (at least) three main challenges to applying model checking
the sorts of open distributed systems we target. First, these systems
must be completed by providing a model of the unknown components. More
importantly, our PAT specifications approximate the behavior of the
underlying actors-- any violations found by model checking these
specifications would have to be validated on the underlying P
implementation of the system. Finally, the state space of these
systems is enormous-- to fully explore the traces would require us to
engineer state (and path) equivalence checkers, frame rules, and other
mechanisms to minimize redundant and uninteresting exploration. This
is precisely why P's existing model checker opts to trade off
engineering complexity for a potentially incomplete search through the
space, driven by a controller. The key insight of this work is that we
can use partial specifications of actors to synthesize controllers
that explore the space of executions in a targeted, property-directed.
As the reviewers note, these controllers are typically incomplete: our
algorithm is only guaranteed to produce a set of controllers that
capture a subset of the feasible executions consistent with provided
specifications (lines 627-629, 776-778). Our experimental evidence
supports our contention that this tradeoff works well in practice.

% It's not immediately apparent to us how we might transform the
% abstract traces that govern the test controller synthesis algorithm to
% a state exploration procedure that a model-checker might use. This is
% because traces under-specify the structure of states. To make such
% traces ``fully explorable'' (Reviewer E) would require us to engineer
% state (and path) equivalence checkers, frame rules, and other
% mechanisms to minimize redundant and uninteresting exploration. We
% propose a lighterweight procedure that leverages traces, not for state
% space exploration, but to search for candidate controllers. In doing
% so, we tradeoff engineering complexity for incompleteness in the
% synthesis algorithm: a synthesized controller is only guaranteed to
% generate some (but not all) feasible executions consistent with
% provided specifications. Our experimental evidence supports our
% contention that this tradeoff works well in practice.

% Three challenges:
%
% 1: Need to close the system
% 2: Complete characterization of components (to be sound)
% 3: Huge state space, need to explore somehow

% QE: In this case we have a complete description of the system (via PAT)
% and that model is executable, meaning that the system should, in
% theory, be fully explorable.


\subsection{Framework}

% Point 3: PATs provide visibility into the system under test,
% allowing for property-directed testing / exploration.

There are two distinguishing features of our framework whose
combination (we believe) separate it from other proposed PBT-style
techniques for distributed systems. First, we target open systems, in
which external clients can inject messages into the SUT. Second, in
contrast to other PBT-style systems that treat the SUT as a black box,
the (local) PAT specifications of the actors provide visibility into
the system in terms of which (global) traces are feasible. This
information enables us to be significantly more targeted in our search
for a violating execution, allowing us to synthesize schedules
\emph{tailored to the property of interest}--- history automata induce
``precondition'' constraints on allowed behaviors that permit the
actor to generate a new event, while prophecy automata induce
``postcondition'' constraints that regulate future actions by the
actor. The capability of dividing symbolic traces into a history and a
future is especially valuable in the concurrent/asynchronous setting
we target.
% Need to additionally emphasize that we are property-directed, PULSE
% and Concuerror appear to be property agnostic-- the former
% implements a "randomized scheduler", while the latter
% "systematically explores process interleaving". Quickstrom, on the
% other hand, does not have information about future events-- it
% randomly chooses an action that is 'enabled' in the current state.

% Quickstrom: No postconditions, not property-directed, actions chosen
% randomly from a set of enabled actions

\textbf{SJ end}

\section{Summary of Proposed Changes}
Concretely, we propose the implement the following changes in the revision, in order to clarify the issues discussed above and the specific criticisms posed by the reviewers addressed below:

\begin{enumerate}
\item We will clarify that \PAT{}s are indeed implemented using SFAs,
  and support both symbolic LTL$_f$ and symbolic regex as frontend
  languages (Reviewers A, B, and E). We will also clarify the
  definition of abstract traces and provide a correctness proof of
  normalization (Reviewer E).

\item To provide a better interpretation of \PAT{}s, we will include
  the type denotation from our supplemental material in the
  revision (Reviewer A).

\item To demonstrate that \PAT{}s enable a high degree of automation,
  we will highlight that the qualifiers in \PAT{} are quantifier-free,
  which guarantees that the VCs are in EPR (Reviewer C). We will also
  provide additional examples of VCs derived from auxiliary type
  judgments (Reviewer A).

\item We will extend our discussion in the evaluation section to
  provide additional details on the testing difficulty of the sorts of
  properties in benchmarks (Reviewer A), elaborate on the
  experimental setup of the second baseline "P+M" (Reviewer A), and
  explain how bugs are injected into our benchmarks (Reviewer B). We
  will also report the average execution time (Reviewer A) and set a
  2-hour time bound (Reviewer B) for the P baseline.

\item We will incorporate all the suggestions made by the reviewers to
  expand our discussion of related work, better contextualizing our
  contributions with respect to other approaches to validating
  distributed systems, including PBT-style approaches (Reviewer D),
  model checking (Reviewer E), and verification (Reviewers B and E).
  Thank you for bringing many of these works to our attention!

\item We will fix all typos and spacing problems identified by the
  reviewers.

\end{enumerate}

\section{Detailed Response}

\paragraph{Reviewer A}
\begin{enumerate}
    \item $\Code{DeriveTerm}$ function. Reviewer is correct where $\Code{DeriveTerm}$ function will specialize the remaining $\globalA A$ into $\epsilon$ to achieve a shorter program. We can synthesize longer programs by forcing the unfolding of $\globalA A$ during synthesis to explore longer abstract traces, although {\sf Clouseau} does not do this for efficiency reasons.
    \item SMT solver usage and VCs. The reviewer is correct; solvers are used in SFA inclusion checks in \textsc{WfHAF} and \textsc{SubHAF}. Following the standard minterm-based SFA algorithm [12], the VCs are proof obligations that all qualifiers of symbolic events (i.e., $\phi$ in $\msgB{op}{\overline{x}}{\phi}$) are satisfiable. We will add VC examples in the revision of our paper.
    \item Execution and assume/assert. The reviewer is correct; the execution needs to ensure that assertions hold, and we do use an SMT solver here.
    \item Data sensitivity. Our properties are data-dependent, so bugs are also data-sensitive, but not in the sense that "a bug is very sensitive to particular values." For example, in the motivating example, we need at least two $\eff{write}$ operations on the same key with different data (as in $A'_\Code{violateRYW}$ on line 357) to lead to a potential inconsistent situation. There are no specific constraints about the values of the key or data here.
    \item More information about "P+M". The "M" is provided from the source of benchmarks, which also contains the human-written "controller" to close the open system. Reviewer is correct that "the best such M would be equivalent to one of your synthesized controllers," but most of the time, M is not, as it is not aware of the property to test.
    \item SubHAF. The reviewer is correct; it is a typo, and the version in the supplemental material is correct. We will include the type denotation from our supplemental material in the revision.
    \item Abstract traces. The reviewer is correct; $A \untilA B$ should be normalized as $\globalA A \seqA B$ (no negation). This is a typo. The meaning of $\globalA \evparenth{\phi}\seqA\Pi$ is exactly as the reviewer states.
\end{enumerate}

\paragraph{Reviewer B}
\begin{enumerate}
    \item Bug injection. The bugs are injected by deleting control flow from the original distributed models directly or by simulating a weakly consistent model (e.g., returning some value written before instead of the last written one).
    \item LTL$_f$ vs. (symbolic) regex. We clarify that \PAT is indeed based on SFAs, and we can accept any front-end language (e.g., symbolic regex, as suggested by reviewers) that can be compiled into SFAs.
\end{enumerate}
We will clarify all these in the paper's revision.

\paragraph{Reviewer C}
\begin{enumerate}
    \item Guarantee constraints are in EPR. We require the qualifiers to be quantifier-free formulas to guarantee that the derived VCs are in EPR. Following the standard minterm-based SFA algorithm [12], the VCs are proof obligations that all qualifiers of symbolic events (i.e., $\phi$ in $\msgB{op}{\overline{x}}{\phi}$) are satisfiable under the type context. The type context is interpreted as prefix universally quantified substitutions (lines $565$ and $580$), which guarantees that the derived VCs are in EPR. We will highlight the constraints on qualifiers in the revision.
    \item Incomplete and getting stuck. The reviewer is correct that our approach is not complete. We would like to clarify that our algorithm is based on backtracking (line 628) and thus will not get stuck when "picking the wrong values or choices." In our motivating example, we show how we shift to another choice when picking the wrong one (lines 753-755).
\end{enumerate}

\paragraph{Reviewer D} Thank you for pointing out the spacing issues in our submission. We will definitely fix them in the revision of our paper!

\paragraph{Reviewer E}

Here is the grammar-checked version of your text:

\begin{enumerate}
    \item Guarantee of witnessing a bug. As mentioned by the reviewer, Theorem 4.5 guarantees that the synthesized controller is type-safe, which provides the bug witness guarantee. Precisely, this guarantee is defined as type soundness (Theorem 3.7), which states "will realize at least one trace consistent with $A$." When the execution is non-deterministic (i.e., the system may randomly pick different executions), we only guarantee that there \emph{exists} one execution that triggers the bug. This is also consistent with the runtime failures observed by reviewer in some benchmarks (e.g., EspressoMachine), as explained in lines 873-876, whose handlers (actors) can non-deterministically fail. The reviewer is correct that runtime failures come from assertion violations, which indicate that the controller's view conflicts with the handlers' view.
    \item Abstract trace and normalization. We clarify that {\sf Clouseau} is indeed based on SFAs, and it can accept any frontend language (e.g., symbolic regex) that can be compiled into SFAs. Our abstract trace can be treated as a symbolic regex without a \emph{top-level} union $\Pi ::= \msgB{op}{\overline{x}}{\phi} ~|~ \Pi \seqA \Pi ~|~ A^*$. Then, we can normalize each SFA into a finite set of abstract traces, as the normalization process preserves the star term instead of unfolding it. The current presentation in our paper attempts to mimic this simple idea in LTL$_f$, which has led to a lot of confusion. We promise to clarify this and provide a correctness proof of normalization in our revised paper.
    \item Liveness. The "liveness property" mentioned in our paper actually refers to "eventually exists within a finite number of steps," following an informal interpretation of liveness: "something good will eventually happen." We apologize for this misuse and promise to clarify it in our paper revision.
    \item Expressivity. We would also like to highlight that \PAT provides constraints over each actor on its \emph{local} view, while the controller has a \emph{global} view. This design ensures that the trace in our approach is not just "one global trace of all events in the system" (Reviewer E), but a merged view of all actors. There can be multiple events for the same operation, coming from different actors under different local views. The global property can then specify whether these events are consistent or not. We provide \PAT examples of "multiple traces per thread" properties (e.g., sequential consistency) and network behaviors expressed in this way, as well as a detailed explanation of \PAT in our benchmarks (e.g., 2PC and RingLeadernominateion) in the attached file of the author response.
    \item Using trace instead of DSL. The DSL is different from "traces together with some assumption(s) and assertion(s)" because of the local variables, which the controller can use to \emph{store} the results from handlers (like register automata). In our case study (lines 961-965), we highlight that our synthesized controller is better than a random controller because it can request a transaction id $\I{tid}$ from the database and use it in future messages. The DSL cannot express this situation without local variables.
    \item Comparison with state-of-the-art baseline. We do compare with the P language, which is indeed a state-of-the-art baseline that has verified realistic distributed models from major cloud vendors in recent years. The reviewer is correct that Mocket is a relevant tool; however, we will not add this comparison in our plan due to the limited revision period.
\end{enumerate}
